\subsubsection{Stable Extrapolation from Controlled Negative Feedback}
\label{sec:controlled_negative_strength}

The causal interventions in Section~\ref{sec:controlled_transmission}
show that far-field negative influence must be controlled. They do not
imply that all negative feedback should be removed. The remaining
question is whether controlled repulsion has a useful regime before
far-field dominance takes over.

We test this by sweeping the effective negative strength relative to
positive attraction, denoted by $q/p$, while keeping the data, labels,
initialization, optimizer, and training horizon fixed. The positive-only
setting corresponds to $q/p=0$. Increasing $q/p$ retains more negative
feedback and increases the repulsive component of the aggregate update.
In the controlled C-U1 geometry, positive attraction anchors the policy
near the observed positive support, while controlled negative repulsion
can push it toward the hidden task optimum; excessive repulsion
over-shoots that useful region.

\begin{figure}[t]
    \centering
    \includegraphics[
        width=0.95\columnwidth
    ]{figures/fig_6_4_1_phase_transition.pdf}

    \vspace{-0.35em}
    \caption{
        \textbf{Stable extrapolation from controlled negative feedback.}
        Increasing the effective negative strength $q/p$ first improves
        held-out-context reward beyond the positive-only ceiling, showing
        that controlled negative feedback can be useful. Larger $q/p$
        then produces over-extrapolation and task-performance collapse.
        Policy shift measures displacement from the positive-only target.
        The color strip reports task collapse, boundary events, and
        NaN/Inf failures separately; each block shows the seed count.
    }
    \label{fig:phase_transition_6_4_1}
\end{figure}

Figure~\ref{fig:phase_transition_6_4_1} shows the resulting transition.
For small and moderate $q/p$, held-out-context reward rises above the
positive-only ceiling. Thus, negative feedback is not merely a source of
instability: when controlled, it can improve beyond the positive-support
solution. As $q/p$ increases further, policy shift continues to grow
while reward falls, indicating over-extrapolation rather than additional
useful generalization. At the largest strengths, task-performance
collapse and boundary events appear, while NaN/Inf failures remain
reported separately.

This strength sweep supplies the missing link between causal rescue and
tapered control. Section~\ref{sec:controlled_transmission} shows that
uncontrolled far-field influence can transmit into collapse. Here, we
show why positive-only learning is not the solution: negative feedback
has a beneficial intermediate regime. Remoteness-aware control should
therefore keep the aggregate update in this useful-repulsion regime while
preventing the far-field tail from dominating.
